
\chapter{Assumptions, Notation, and Proof Discipline}
\label{ch:assumptions}

This chapter establishes the formal scaffolding for Part~III.  It fixes
notation that is stable across all theorems, catalogues the eight
standing assumptions, defines the proof dependency matrix, and sets the
disciplinary standards that every subsequent proof must satisfy.  The
flagship DC3S ladder in Chapters~\ref{ch:theorem-oasg}--\ref{ch:temporal-extensions}
contains eight principal theorems (T1--T8), but it is not the entire theorem
surface of the dissertation.  Part~I establishes two foundational precursor
theorems (S1: illusion existence under dropout; S2: informal DC3S guarantee);
the flagship chapters add 5~lemmas, 10~corollaries, and 7~propositions; and
16~standing assumptions ground every formal claim.  The integrated surface
totals \textbf{67 formal items} (18~theorems, 5~lemmas, 7~propositions,
10~corollaries, 11~definitions, 16~assumptions), all audited in
Appendix~\ref{app:integrated-theorem-surface-register} and summarised in
Table~\ref{tab:integrated-theorem-surface-summary} at the end of this chapter.

\section{Stable Notation}

Throughout Part~III, the following symbols have fixed meaning:

\begin{center}
\begin{tabular}{cl}
\toprule
\textbf{Symbol} & \textbf{Meaning} \\
\midrule
$x_t$ & True (plant-internal) state at step $t$ \\
$\hat{x}_t$ & Observed (telemetry-delivered) state at step $t$ \\
$a_t$ & Action (dispatch decision) at step $t$ \\
$a_t^{\mathrm{cand}}$ & Candidate action proposed from the observation history \\
$a_t^{\mathrm{safe}}$ & Repaired action returned by the ORIUS kernel \\
$U_t$ & Observation-consistent state set built from degraded telemetry \\
$\mathcal{S}$ & Safe state set: $\mathcal{S} = \{x : \phi(x) \ge 0\}$ \\
$\phi(x_t)$ & Safety predicate: $\phi(x_t) = 1$ iff $x_t$ satisfies all constraints \\
$w_t$ & Reliability score, $w_t \in [0,1]$ \\
$\alpha$ & Conformal miscoverage level (nominal $\alpha = 0.10$) \\
$q_t$ & Conformal quantile (half-width of the prediction interval) \\
$\mathcal{A}_t$ & Tightened safe-action set at step $t$ \\
$\mathcal{H}_t$ & Causal observation history available at time $t$ \\
$T$ & Episode horizon (number of steps) \\
$V$ & Total violation count: $V = \sum_{t=1}^T \mathbf{1}[\lnot\,\phi(x_t)]$ \\
$\bar{w}$ & Episode-average reliability: $\bar{w} = T^{-1}\sum_t w_t$ \\
\bottomrule
\end{tabular}
\end{center}

\paragraph{Battery instantiation.}
In the battery domain, $x_t = \text{SOC}_t$ (MWh),
$a_t = (P_t^{\text{chg}}, P_t^{\text{dis}})$ (MW), and
$\phi(x_t) = \mathbf{1}[\text{SOC}_{\min} \le x_t \le \text{SOC}_{\max}]$.

\section{Assumption IDs (A1--A8)}

Each assumption is assigned a permanent identifier and stated in both
domain-general and battery-specific form.

\begin{assumption}[A1: Bounded model error]
\label{assm:A1}
The one-step prediction error $\|x_{t+1} - f(x_t, a_t)\|$ is bounded
by a known constant $\epsilon_{\text{model}}$.
\emph{Battery:} $|x_{t+1} - (x_t + \eta\,P_t\,\Delta t)| \le \epsilon_{\text{model}}$.
\end{assumption}

\begin{assumption}[A2: Bounded telemetry error]
\label{assm:A2}
The observation error $\|\hat{x}_t - x_t\|$ is bounded by a
reliability-dependent function: $\|\hat{x}_t - x_t\| \le g(w_t)$,
where $g$ is monotone decreasing.
\end{assumption}

\begin{assumption}[A3: Feasible safe repair]
\label{assm:A3}
For every state $x_t$ satisfying $\phi(x_t) = 1$, there exists at
least one action $a_t \in \mathcal{A}_t$ such that $\phi(x_{t+1}) = 1$.
\end{assumption}

\begin{assumption}[A4: Known dynamics]
\label{assm:A4}
The transition function $f(x_t, a_t)$ is known to the controller up
to the model error bounded in A1.
\emph{Battery:} the SOC transition $x_{t+1} = x_t + \eta\,P_t\,\Delta t$
is known with efficiency $\eta$ identified offline.
\end{assumption}

\begin{assumption}[A5: Monotone bounded uncertainty inflation]
\label{assm:A5}
The inflation function $\iota(w_t) = q_t / (w_t + \epsilon)$ is
monotone decreasing in $w_t$ and bounded: $\iota(w_t) \le \bar{\iota}$
for all $w_t \in [0,1]$.
\end{assumption}

\begin{assumption}[A5$'$: Calibration regularity]
\label{assm:A5prime}
The telemetry error bound $g(w_t)$ from Assumption~A2 satisfies
$g(1) = 0$ and is Lipschitz near $w_t = 1$: there exist constants
$c_g > 0$ and $\delta_w \in (0,1)$ such that
$g(w_t) \le c_g\,(1-w_t)$ for all $w_t \in [1-\delta_w, 1]$.
This calibration regularity aligns the error bound with the reliability
score and is used in the information-theoretic lower bound (T10).
\end{assumption}

\begin{assumption}[A6: Bounded detector lag]
\label{assm:A6}
The OQE reliability score $w_t$ reflects a telemetry fault within at
most $\ell$ steps of fault onset: if a fault begins at step $t_f$,
then $w_{t_f + \ell} \le w_{\text{alert}}$.
\end{assumption}

\begin{assumption}[A6$'$: $\varphi$-mixing fault process]
\label{assm:A6prime}
The fault injection process $\{F_t\}$ is $\varphi$-mixing with
geometric decay: there exist constants $C_\phi > 0$ and $\rho\in(0,1)$
such that the $\varphi$-mixing coefficient satisfies $\varphi(k) \le
C_\phi\,\rho^k$ for all $k \ge 0$.  A buffer of
$\tau_{\mathrm{mix}} = \lceil\log(1/\varepsilon)/|\log\rho|\rceil$
steps makes any two windows $\varepsilon$-approximately independent in
total variation.  Assumption~A6$'$ is used in the universal
impossibility proof (T9) and is satisfied by any i.i.d.\ or
geometrically ergodic Markov fault model.
\end{assumption}

\begin{assumption}[A7: Causal certificate update]
\label{assm:A7}
The certificate $\mathcal{C}_t$ depends only on information available
at or before step $t$: $\mathcal{C}_t = h(\hat{x}_{1:t}, a_{1:t-1})$.
\end{assumption}

\begin{assumption}[A8: Admissible fallback policy]
\label{assm:A8}
There exists a fallback policy $\pi_{\text{fb}}$ that, starting from any
safe state, maintains $\phi(x_t) = 1$ for all future steps under worst-case
bounded disturbance.
\emph{Battery:} the zero-dispatch policy ($P_t = 0$) is always admissible
because SOC is unchanged (up to model error).
\end{assumption}

\paragraph{Executable universal kernel.}
The canonical code surface for these assumptions now lives in
\path{src/orius/universal_theory/}.  That package exposes typed runtime
objects for observation packets, reliability assessments,
observation-consistent state sets, safe-action sets, repair decisions, and
safety certificates.  The battery-domain proofs remain the reference
instantiation, but the reusable universal claim now attaches to this typed
kernel rather than to an untyped adapter narrative.

\IfFileExists{reports/publication/assumption_code_traceability.tex}{%
  \begin{table*}[htbp]
\centering\footnotesize
\caption{Assumption-to-code traceability for the universal degraded-observation kernel.}
\label{tab:assumption-code-traceability}
\begin{tabular}{p{0.8cm}p{3.2cm}p{6.1cm}p{3.0cm}}
\toprule
\textbf{ID} & \textbf{Assumption} & \textbf{Canonical code surface} & \textbf{Runtime status} \\
\midrule
A1 & Bounded model error & Domain instantiation and battery reference helpers in \path{src/orius/universal_theory/battery_instantiation.py} & Declared by model/artifact surface \\
A2 & Bounded telemetry error & \path{compute_oqe()}, \path{build_reliability_assessment()}, observation-consistent state construction & Checked through reliability and state envelopes \\
A3 & Feasible safe repair & \path{tighten_action_set()}, \path{repair_action()}, \path{build_repair_decision()} & Checked per step \\
A4 & Known dynamics & \path{DomainInstantiation} contract and domain-specific transition models & Declared by domain instantiation \\
A5 & Monotone bounded inflation & \path{verify_inflation_geq_one()} and clipped inflation in \path{build_observation_consistent_state_set()} & Checked per step \\
A6 & Bounded detector lag & OQE/drift artifacts and locked replay audits & Artifact-validated; not fully checked at runtime \\
A7 & Causal certificate update & \path{execute_universal_step()} and \path{SafetyCertificate.from_payload()} & Checked structurally \\
A8 & Admissible fallback policy & \path{SafetySpec.fallback_action}, battery zero-dispatch fallback, fallback branch in \path{build_repair_decision()} & Checked when fallback exists \\
\bottomrule
\end{tabular}
\end{table*}
%
}{%
  \begin{table*}[htbp]
\centering\footnotesize
\caption{Assumption-to-code traceability for the universal degraded-observation kernel.}
\label{tab:assumption-code-traceability}
\begin{tabular}{p{0.8cm}p{3.2cm}p{6.1cm}p{3.0cm}}
\toprule
\textbf{ID} & \textbf{Assumption} & \textbf{Canonical code surface} & \textbf{Runtime status} \\
\midrule
A1 & Bounded model error & Domain instantiation and battery reference helpers in \path{src/orius/universal_theory/battery_instantiation.py} & Declared by model/artifact surface \\
A2 & Bounded telemetry error & \path{compute_oqe()}, \path{build_reliability_assessment()}, observation-consistent state construction & Checked through reliability and state envelopes \\
A3 & Feasible safe repair & \path{tighten_action_set()}, \path{repair_action()}, \path{build_repair_decision()} & Checked per step \\
A4 & Known dynamics & \path{DomainInstantiation} contract and domain-specific transition models & Declared by domain instantiation \\
A5 & Monotone bounded inflation & \path{verify_inflation_geq_one()} and clipped inflation in \path{build_observation_consistent_state_set()} & Checked per step \\
A6 & Bounded detector lag & OQE/drift artifacts and locked replay audits & Artifact-validated; not fully checked at runtime \\
A7 & Causal certificate update & \path{execute_universal_step()} and \path{SafetyCertificate.from_payload()} & Checked structurally \\
A8 & Admissible fallback policy & \path{SafetySpec.fallback_action}, battery zero-dispatch fallback, fallback branch in \path{build_repair_decision()} & Checked when fallback exists \\
\bottomrule
\end{tabular}
\end{table*}
%
}

\section{Proof Dependency Matrix}

Table~\ref{tab:proof-deps} shows which assumptions each theorem requires.

\begin{table}[ht]
\centering
\caption{Assumption dependencies of main theorems.}
\label{tab:proof-deps}
\begin{tabular}{lccccccccc}
\toprule
\textbf{Theorem} & \textbf{A1} & \textbf{A2} & \textbf{A3} & \textbf{A4}
  & \textbf{A5} & \textbf{A6} & \textbf{A7} & \textbf{A8} \\
\midrule
OASG Existence       & \checkmark & \checkmark &            & \checkmark &            &            &            &            \\
Safety Preservation  & \checkmark & \checkmark & \checkmark & \checkmark & \checkmark &            & \checkmark &            \\
Core Bound           & \checkmark & \checkmark &            & \checkmark & \checkmark & \checkmark & \checkmark &            \\
No Free Safety       &            & \checkmark &            &            &            &            &            &            \\
Certificate Horizon  & \checkmark & \checkmark &            & \checkmark & \checkmark & \checkmark & \checkmark &            \\
Graceful Degradation & \checkmark & \checkmark & \checkmark & \checkmark & \checkmark &            &            & \checkmark \\
\bottomrule
\end{tabular}
\end{table}

\subsection{Extended Monograph Theorem Program}

The flagship dependency matrix above covers the battery proof ladder.  The
monograph-level extension adds T9--T11, each of which answers a distinct
control-theory question and introduces an explicit new obligation.

\begin{table}[htbp]
\centering
\caption{Extended theorem program beyond the battery ladder.}
\label{tab:extended-theorem-program}
\begin{tabular}{p{1.3cm}p{4.4cm}p{3.2cm}p{3.7cm}}
\toprule
\textbf{Theorem} & \textbf{Question answered} & \textbf{New assumptions} & \textbf{Role} \\
\midrule
T9 & Is reliability-aware repair structurally necessary under persistent degradation? & A6$'$ + boundary reachability from A4 & Necessity theorem for degraded-observation control \\
\addlinespace
T10 & Can the $\alpha(1-\bar{w})T$ envelope be improved by a different controller alone? & A5$'$ + conformal coverage condition & Lower-bound scaling theorem \\
\addlinespace
T11 & Which structural obligations let the battery argument transfer to a new domain? & Explicit soundness, repair-membership, and fallback obligations & Template transfer theorem \\
\bottomrule
\end{tabular}
\end{table}

\section{Marginal Versus Conditional Coverage}

A critical distinction throughout Part~III:
\begin{description}
  \item[Marginal coverage.] $\Pr[\,x_t \in \hat{C}_t\,] \ge 1 - \alpha$
    averaged over the randomness of calibration and test points.
    This is what CQR provides.
  \item[Conditional coverage.] $\Pr[\,x_t \in \hat{C}_t \mid \hat{x}_t\,]
    \ge 1 - \alpha$ for every observed value.  Conditional coverage is
    strictly stronger and is \emph{not} guaranteed by CQR without
    additional distributional assumptions.
\end{description}

The ORIUS theorems are stated under marginal coverage.  Where conditional
coverage is needed (e.g., for the per-step risk lemma), we explicitly
invoke the monotone inflation assumption~(A5) as a surrogate.

\section{One-Step Versus Multi-Step Claims}

The theorems in Chapters~\ref{ch:theorem-oasg}--\ref{ch:core-bound}
operate on a \emph{per-step} basis: each step's safety is analysed
conditionally on the current state.  The core bound
(Chapter~\ref{ch:core-bound}) aggregates per-step risks into an
episode-level expectation via linearity of expectation.

Multi-step claims (certificate horizon, graceful degradation) require
additional assumptions about temporal dependence; these are stated
explicitly in Chapter~\ref{ch:temporal-extensions}.

\section{The Theorem Ladder}

The theorems form a logical ladder:

\begin{enumerate}
  \item \textbf{OASG Existence} (Chapter~\ref{ch:theorem-oasg}):
    The problem is real---observed safety does not imply true safety.
  \item \textbf{Safety Preservation} (Chapter~\ref{ch:theorem-safety}):
    The shield solves the problem---true safety is restored.
  \item \textbf{Core Bound} (Chapter~\ref{ch:core-bound}):
    The solution has a quantifiable cost---expected violations are bounded.
  \item \textbf{No Free Safety} (Chapter~\ref{ch:no-free-safety}):
    The cost is unavoidable---ignoring observation quality breaks safety.
  \item \textbf{Temporal Extensions} (Chapter~\ref{ch:temporal-extensions}):
    The guarantee extends across time---certificates expire gracefully.
\end{enumerate}

Each rung depends on the ones below it, and no rung can be removed
without breaking the chain.

\section{Theorem-Writing Discipline}

Every theorem in this dissertation follows a fixed template:

\begin{enumerate}
  \item \textbf{Statement}: formal claim with quantifiers.
  \item \textbf{Assumptions used}: list of A-IDs.
  \item \textbf{Proof}: complete derivation.
  \item \textbf{Artifact link}: pointer to the empirical evidence that
    instantiates the theorem in the battery domain.
  \item \textbf{Scope disclaimer}: ``What this theorem does not prove.''
\end{enumerate}

\section{Common Theorem Failure Modes}

To inoculate against reviewer objections, we catalogue failure modes
that a careless proof might exhibit:

\begin{enumerate}
  \item \emph{Vacuous truth.}  If the assumption set is never satisfied,
    the theorem is true but useless.  We verify non-vacuousness by
    exhibiting the battery instantiation.
  \item \emph{Circular dependence.}  Theorem~$X$ assumes the conclusion
    of Theorem~$Y$, and vice versa.  The dependency matrix
    (Table~\ref{tab:proof-deps}) prevents this.
  \item \emph{Hidden conditional.}  A marginal claim is stated as if it
    were conditional.  We use the marginal/conditional distinction
    explicitly.
  \item \emph{Infinite-sample limit.}  A result that holds only
    asymptotically is presented as finite-sample.  All ORIUS bounds
    are stated for finite~$T$.
\end{enumerate}

\section{Integrated Theorem Surface Summary}

Table~\ref{tab:integrated-theorem-surface-summary} enumerates the full
formal surface of the dissertation.  The ten unique theorems (S1, S2,
T1--T8) and their appendix restatements (C.1--C.9) all pass the
integrated theorem gate: source file exists, code anchor verified,
artifact locked, and appendix mapping confirmed.

\IfFileExists{reports/publication/theorem_surface_summary.tex}{%
  \begin{table}[htbp]
\centering
\caption{YAML-canonical theorem surface summary.}
\begin{tabular}{lr}
\toprule
Environment & Count \\
\midrule
theorem & 27 \\
lemma & 8 \\
proposition & 0 \\
corollary & 2 \\
definition & 0 \\
assumption & 14 \\
remark & 0 \\
\midrule
total & 51 \\\\
\bottomrule
\end{tabular}
\end{table}
%
}{%
  \begin{table}[htbp]
\centering
\caption{YAML-canonical theorem surface summary.}
\begin{tabular}{lr}
\toprule
Environment & Count \\
\midrule
theorem & 27 \\
lemma & 8 \\
proposition & 0 \\
corollary & 2 \\
definition & 0 \\
assumption & 14 \\
remark & 0 \\
\midrule
total & 51 \\\\
\bottomrule
\end{tabular}
\end{table}
%
}

\section{Chapter Summary}

This chapter provides the stable notational, assumptive, and
methodological foundation for every proof in Part~III.  The eight
assumptions (A1--A8) are individually justified in both domain-general
and battery-specific terms, the proof dependency matrix prevents
circular reasoning, and the theorem-writing discipline ensures that
every claim is bounded by explicit scope disclaimers.  The integrated
theorem surface of 67 formal items---spanning 18 theorems, 5 lemmas,
7 propositions, 10 corollaries, 11 definitions, and 16 assumptions---is
the formal depth this dissertation contributes to the CPS safety
literature.
